\chapter{Présentation du thème du stage}

\section{Champ de l'étude et objet du stage}
Le thème du stage est la réalisation de la couche \textbf{frontend} d'une plateforme éducative moderne, en \textbf{React + TypeScript}, avec un objectif d'ergonomie, de clarté fonctionnelle et d'intégration fluide avec les services backend.

L'objet de mon lot couvre :
\begin{itemize}
	\item l'architecture des pages et la navigation ;
	\item les interfaces des espaces enseignant et étudiant ;
	\item l'intégration des appels API ;
	\item la validation fonctionnelle des scénarios utilisateur.
\end{itemize}

\section{Planning prévisionnel du stage}
\begin{longtable}{p{0.25\textwidth}p{0.65\textwidth}}
\toprule
Semaine 1 & Analyse des besoins UI/UX et définition des parcours utilisateurs \\
Semaine 2 & Mise en place de l'architecture frontend et du routage protégé \\
Semaine 3 & Développement des pages enseignant et intégration API associée \\
Semaine 4 & Développement des pages étudiant, validation des flux et finalisation UI \\
\bottomrule
\end{longtable}
